\subsection{多固有値Ore領域の非Jordan基底の認識}
\label{subsec:cubic-spectral-Ore-reduction}
\begin{lemma}[三次格子の中心を通る直線]\label{lem:cubic-lattice-line}
$k,m,n\ge0$、$k+m+n=3$ とし、$(s,t)\ne(0,0)$ を固定する。
$(m-1)s+(n-1)t=0$ を満たす全格子点では、三座標のうち少なくとも一つが
共通して正である。
\end{lemma}
\begin{proof}
$(1,1,1)$ 以外の格子点の方向 $(m-1,n-1)$ を符号と共通因子を除いて列挙すると
\[
 (1,1),(1,0),(1,-1),(1,-2),(0,1),(2,-1)
\]
だけである。それぞれの直線上の点で共通して正な座標は順に $k,n,k,n,m,m$。
いずれの方向も該当しないときは $(1,1,1)$ のみである。
\end{proof}

\begin{proposition}[可換平面上の非半単純可換作用の排除]
\label{prop:cubic-commuting-Ore-no-shear}
$T=k[x,y]$、$f,g\in\operatorname{GL}(T_1)$、$fg=gf$、$q\ne0$ とし、
$B=T[b;f][c;\widetilde g]$、$\widetilde g(b)=q^{-1}b$ とする。
$f,g$ の少なくとも一つが非半単純なら、$B$ はH3包絡にならない。
従ってこの可換基底の全H3部分はskew包絡に属する。
\end{proposition}
\begin{proof}
可換な二つの可逆行列の一方が非半単純なら、共通基底で
\[
 f(x)=\alpha x,\ f(y)=\alpha(y+s x),\qquad
 g(x)=\beta x,\ g(y)=\beta(y+t x),\qquad(s,t)\ne(0,0)
\]
と書ける。Ore中山公式は
$\mu_B|_T=(fg)^{-1}$、$\mu_B(b)\in k^*b$、$\mu_B(c)\in k^*c$ を与える。
H3の必要正常三次をPBW基底で
$G=\sum_{m+n\le3}F_{mn}(x,y)b^mc^n$ と書く。
$b,c$ に対する正規式から、各 $F_{mn}$ は $f,g$ の固有ベクトルである。
少なくとも一つのshearが非零なので、次数 $k=3-m-n$ のその固有ベクトルは $kx^k$ のみ。
従って $G$ は $x^kb^mc^n$ の線型結合となる。

$Gx=\mu_B^{-1}(x)G=\alpha\beta xG$ のPBW係数比較により、
非零項ごとに $\alpha^m\beta^n=\alpha\beta$。
$Gy=\alpha\beta(y+(s+t)x)G$ と比較すると
\[
 (m-1)s+(n-1)t=0
\]
を得る。前補題から全非零項に共通する $x,b,c$ のいずれかが存在する。
これら三生成元は互いにskew可換なので、$G$ はその一次生成元を左因子として持つ。
三次商ではこの一次因子と残る非零二次因子が零因子となり、整域性に反する。
一方、$f,g$ がともに半単純なら可換性から同時対角化でき、$B$ はskewである。
\end{proof}

\begin{proposition}[反交換基底の反対角作用の排除]
\label{prop:cubic-anticommuting-Ore-no-swap}
定理\ref{thm:cubic-211-nonresonant-Ore}の領域で、基底平面が
$T=k\langle x,y\rangle/(xy+yx)$ なら、全H3包絡はskewである。
\end{proposition}
\begin{proof}
$T$ の自己同型は対角または反対角。
一方だけが反対角なら、可換性から他方はscalarであり、$fg$ は反対角。
$\mu_T=-I$ とOre公式から $\mu_B|_U=-(fg)^{-1}$ は二つの異なる固有値を持ち、
重複2の半単純固有空間という仮定に反する。
両方対角ならskewである。
両方反対角なら可換性から $fg=dI$、$d\ne0$。

正常三次のPBW係数 $F_{mn}\in T_{3-m-n}$ に対する $T$ との正常式は
\[
 F_{mn}f^mg^n(u)=-d\,uF_{mn}\quad(u=x,y)
\]
となる。$m+n$ が奇数なら $f^mg^n$ は反対角である。
非零 $F_{mn}$ が存在すると、そのnormalizerは反対角自己同型のscalar倍となる。
しかし正常元のnormalizerは中心 $x^2,y^2$ を点ごとに固定する必要があり、
反対角作用はこれらを交換するため不可能。
$m+n$ が偶数なら $f^mg^n$ はscalarであり、$F_{mn}$ の次数は奇数である。
$F x=\rho xF$、$F y=\rho yF$ を満たす奇数次の元は零である。
実際、PBW単項式 $x^iy^j$ の二つの符号は $(-1)^j,(-1)^i$ で、
$i+j$ が奇数なら等しくならない。従って全係数が零となり $W_1(B)=0$。
\end{proof}

\begin{corollary}[二次非共鳴 $(2,1,1)$ の残余]\label{cor:cubic-211-Jordan-residual}
定理\ref{thm:cubic-211-nonresonant-Ore}のH3包絡は、既分類のskew包絡、
または基底 $T$ がJordan平面である逐次Ore拡大に限られる。
\end{corollary}
\begin{proof}
二元AS2平面は量子平面またはJordan平面である。
可換平面と反交換平面は上の二命題で処理した。
残る量子平面では全自己同型が対角であり、全包絡がskewとなる。
Jordan残余の存在や許容性はここでは結論しない。
\end{proof}
